%% toc: P15 165 165.5
\section{Ranking theory: contraction, independence and consonant belief functions}
\label{sec:p15-w02b-p15-039-annex}

\begin{annexitem}{annex:p15-w02b-p15-039-levi}{thm:p15-w02b-p15-039-levi-harper}{Levi and Harper identities}
\begin{proof}
(i) Let \(\lambda\) be regular with world ranks \(r_\lambda\).  By
Theorem~\ref{thm:p15-rat53-rank-laws},
\(\lambda(\neg A)=\min_{w\in\neg A}r_\lambda(w)\), with
\(\min\varnothing=\infty\).  Hence \(\lambda(\neg A)>0\) exactly when no
world of \(\neg A\) has rank zero, that is, exactly when
\(M_0(\lambda)\subseteq A\).  For the expansion, if
\(A\supseteq B\land E\) with \(B\in\mathcal B_0(\lambda)\), then
\(B\supseteq M_0(\lambda)\) and so \(A\supseteq M_0(\lambda)\land E\).
Conversely \(M_0(\lambda)\in\mathcal B_0(\lambda)\) by the first statement,
so every \(A\supseteq M_0(\lambda)\land E\) lies in
\(\mathcal B_0(\lambda)+E\), witnessed by \(B=M_0(\lambda)\).

Since \(\kappa\) is regular and \(\varnothing\ne E\ne W\), both
\(\kappa(E)\) and \(\kappa(\neg E)\) are finite.  By
Theorem~\ref{thm:p15-rat53-update}, \(\kappa^{a,b}_E\) has world ranks
\(r(w)-\kappa(E)+a\) on \(E\) and \(r(w)-\kappa(\neg E)+b\) on
\(\neg E\).  These are finite, so every such update of a regular ranking is
regular.  On \(E\), \(r(w)-\kappa(E)\geq0\) with equality exactly on
\(\operatorname{Min}_\kappa(E)\), and likewise on \(\neg E\).  With
\((a,b)=(0,n)\) and \(n\geq1\), every world of \(\neg E\) has new rank at
least \(n\geq1\), so the zero set is \(\operatorname{Min}_\kappa(E)\).  With
\((a,b)=(0,0)\) the zero set is
\(\operatorname{Min}_\kappa(E)\lor\operatorname{Min}_\kappa(\neg E)\).

Two facts are used below, for any proposition \(X\) with
\(\varnothing\ne X\ne W\).
(F1) If \(\kappa(X)=0\), then
\(\operatorname{Min}_\kappa(X)=M_0(\kappa)\land X\), because both consist
of the worlds of \(X\) with rank zero.
(F2) If \(\kappa(X)>0\), then no world of rank zero lies in \(X\), so
\(M_0(\kappa)\subseteq\neg X\); the negation law gives
\(\kappa(\neg X)=0\), and then
\(\operatorname{Min}_\kappa(\neg X)=M_0(\kappa)\).

(ii) If \(\kappa(\neg E)=0\), then \(\kappa^{-}_E=\kappa\) is regular, the
inclusion is an equality, and \(E\notin\mathcal B_0(\kappa)\) by the
definition of \(\mathcal B_0\).  If \(\kappa(\neg E)>0\), then
\(\kappa^{-}_E=\kappa^{0,0}_E\) is regular by~(i).  By (F2) with
\(X=\neg E\), \(\operatorname{Min}_\kappa(E)=M_0(\kappa)\), so by~(i)
\[
 M_0(\kappa^{-}_E)=M_0(\kappa)\lor\operatorname{Min}_\kappa(\neg E)
 \supseteq M_0(\kappa).
\]
Every \(A\) containing \(M_0(\kappa^{-}_E)\) therefore contains
\(M_0(\kappa)\), which is the inclusion.  The set
\(\operatorname{Min}_\kappa(\neg E)\) is nonempty, because \(\neg E\) is
nonempty and finite and its minimum rank is attained, and it is disjoint
from \(E\).  Hence \(M_0(\kappa^{-}_E)\not\subseteq E\), and
\(E\notin\mathcal B_0(\kappa^{-}_E)\) by~(i).

(iii) By~(i) the left side is
\(\{A:\operatorname{Min}_\kappa(E)\subseteq A\}\).  The case condition for
\(\kappa^{-}_{\neg E}\) is the value of \(\kappa(\neg\neg E)=\kappa(E)\).
If \(\kappa(E)=0\), then \(\kappa^{-}_{\neg E}=\kappa\), and by~(i) and
(F1) the right side is
\(\{A:M_0(\kappa)\land E\subseteq A\}=\{A:\operatorname{Min}_\kappa(E)\subseteq A\}\).
If \(\kappa(E)>0\), then \(\kappa^{-}_{\neg E}=\kappa^{0,0}_{\neg E}\).  Part~(i)
applied to \(\neg E\), which also satisfies \(\varnothing\ne\neg E\ne W\),
gives its zero set
\(\operatorname{Min}_\kappa(\neg E)\lor\operatorname{Min}_\kappa(E)\), which
by (F2) with \(X=E\) is \(M_0(\kappa)\lor\operatorname{Min}_\kappa(E)\).  Its
intersection with \(E\) is \(\operatorname{Min}_\kappa(E)\), since
\(M_0(\kappa)\subseteq\neg E\) by (F2) and
\(\operatorname{Min}_\kappa(E)\subseteq E\).  By~(i) the right side is again
\(\{A:\operatorname{Min}_\kappa(E)\subseteq A\}\).

(iv) By~(i) applied to \(\neg E\),
\(\mathcal B_0(\kappa^{0,n}_{\neg E})=\{A:\operatorname{Min}_\kappa(\neg E)\subseteq A\}\),
so the right side is
\(\{A:M_0(\kappa)\lor\operatorname{Min}_\kappa(\neg E)\subseteq A\}\).
If \(\kappa(\neg E)=0\), then (F1) with \(X=\neg E\) gives
\(\operatorname{Min}_\kappa(\neg E)=M_0(\kappa)\land\neg E\subseteq M_0(\kappa)\),
so the right side is \(\mathcal B_0(\kappa)=\mathcal B_0(\kappa^{-}_E)\).
If \(\kappa(\neg E)>0\), the display in the proof of~(ii) gives
\(M_0(\kappa^{-}_E)=M_0(\kappa)\lor\operatorname{Min}_\kappa(\neg E)\), and
(i) gives equality with the right side.

(v) If \(\kappa(E)=0\), then \(\kappa^{-}_{\neg E}=\kappa\) and there is
nothing to prove.  If \(\kappa(E)>0\), put
\(\lambda=\kappa^{0,0}_{\neg E}\).  By Theorem~\ref{thm:p15-rat53-update},
\(\lambda(E)=\lambda(\neg E)=0\) and \(\lambda\) has the old conditional
rank functions within both cells:
\(\lambda(A\mid E)=\kappa(A\mid E)\) and
\(\lambda(A\mid\neg E)=\kappa(A\mid\neg E)\) for every \(A\).  The defining
formula of Definition~\ref{def:p15-rat53-soft-update} then gives, for every
\(A\),
\[
\begin{aligned}
 \lambda^{0,n}_E(A)&=\min\{\lambda(A\mid E),\ \lambda(A\mid\neg E)+n\}\\
 &=\min\{\kappa(A\mid E),\ \kappa(A\mid\neg E)+n\}=\kappa^{0,n}_E(A).
\end{aligned}
\]

For Example~\ref{def:p15-w02b-p15-039-levi-example}: with
\(\kappa_i(q)=0\) and \(\kappa_i(\neg q)=1\), the world formula of~(i) for
\((\kappa_i)^{0,0}_q\) leaves the \(q\)-worlds \((1,1),(0,1)\) at \(0,3\)
and lowers \((1,0),(0,0)\) by one: to \(0,1\) for \(\kappa_1\) and to
\(1,0\) for \(\kappa_2\).  The zero sets are \(\{(1,1),(1,0)\}\) and
\(\{(1,1),(0,0)\}\), the truth sets of \(p\) and of \(p\leftrightarrow q\).
For \(\kappa_1\), \(\operatorname{Min}_{\kappa_1}(\neg q)=\{(1,0)\}\), which
is the intersection of the first zero set with the truth set of \(\neg q\),
and \(M_0(\kappa_1)\lor\{(1,0)\}\) is the first zero set, as~(iii)
and~(iv) require.
\end{proof}
\end{annexitem}

\begin{annexitem}{annex:p15-w02b-p15-039-independence}{thm:p15-w02b-p15-039-independence}{Ranking independence}
\begin{proof}
All four cells are finite, so the marginals \(\kappa(A)=\min\{a,b\}\),
\(\kappa(\neg A)=\min\{c,d\}\), \(\kappa(B)=\min\{a,c\}\) and
\(\kappa(\neg B)=\min\{b,d\}\) are finite, by the binary minimum law of
Definition~\ref{def:p15-rat53-ranking}.  By
Definition~\ref{def:p15-rat53-conditionals},
\(\kappa(B\mid A)=a-\kappa(A)\), \(\kappa(\neg B\mid A)=b-\kappa(A)\),
\(\kappa(B\mid\neg A)=c-\kappa(\neg A)\) and
\(\kappa(\neg B\mid\neg A)=d-\kappa(\neg A)\), all finite; hence
\(\tau(B\mid A)=b-a\) and \(\tau(B\mid\neg A)=d-c\).

\emph{(1)\(\Leftrightarrow\)(2).}  Both condition cells have finite rank,
so the reason relation of Definition~\ref{def:p15-rat53-conditionals}
applies to \(A\) and to \(\neg A\).  \(A\) is a reason for \(B\) exactly
when \(\tau(B\mid A)>\tau(B\mid\neg A)\), and \(\neg A\) is a reason for
\(B\) exactly when \(\tau(B\mid\neg A)>\tau(B\mid\neg\neg A)=\tau(B\mid A)\).
The two values are finite integers, so neither inequality holds exactly when
they are equal.

\emph{(2)\(\Leftrightarrow\)(3).}  \(b-a=d-c\) exactly when \(a+d=b+c\).

\emph{(3)\(\Rightarrow\)(4).}  Put \(\delta=b-a=d-c\).  Suppose first
\(\delta\geq0\).  Then \(b\geq a\) and \(d\geq c\), so \(\kappa(A)=a\) and
\(\kappa(\neg A)=c\), and \(\kappa(\neg B)=\min\{a+\delta,c+\delta\}
=\kappa(B)+\delta\).  By the negation law of
Theorem~\ref{thm:p15-rat53-rank-laws}, one of \(\kappa(B)\) and
\(\kappa(\neg B)\) is zero; if \(\kappa(B)>0\), then
\(\kappa(\neg B)=\kappa(B)+\delta>0\) as well, which is impossible.  So
\(\kappa(B)=0\) and \(\kappa(\neg B)=\delta\).  Now
\(a=\kappa(A)+\kappa(B)\), \(b=a+\delta=\kappa(A)+\kappa(\neg B)\),
\(c=\kappa(\neg A)+\kappa(B)\) and \(d=c+\delta=\kappa(\neg A)+\kappa(\neg B)\),
which are the four equations of~(4).  If \(\delta<0\), exchange the roles of
\(B\) and \(\neg B\).  The exchange swaps \(a\) with \(b\) and \(c\) with
\(d\); it leaves the equation \(a+d=b+c\) and the set of four equations
in~(4) unchanged, and it replaces \(\delta\) by \(a-b=-\delta>0\).  The
first case then gives~(4).

\emph{(4)\(\Rightarrow\)(3).}  Adding the equations for the cells
\(A\land B\) and \(\neg A\land\neg B\) gives
\(a+d=\kappa(A)+\kappa(\neg A)+\kappa(B)+\kappa(\neg B)\), and adding those
for \(A\land\neg B\) and \(\neg A\land B\) gives the same sum for \(b+c\).

\emph{(4)\(\Rightarrow\)(5).}  By~(4),
\(\kappa(B\mid A)=\kappa(A\land B)-\kappa(A)=\kappa(B)\), and in the same
way \(\kappa(B\mid\neg A)=\kappa(B)\) and
\(\kappa(\neg B\mid A)=\kappa(\neg B\mid\neg A)=\kappa(\neg B)\).  The
update is defined because \(\kappa(A),\kappa(\neg A)<\infty\), and by
Definition~\ref{def:p15-rat53-soft-update}
\[
 \kappa^{x,y}_A(B)=\min\{\kappa(B\mid A)+x,\ \kappa(B\mid\neg A)+y\}
 =\kappa(B)+\min\{x,y\}=\kappa(B).
\]
The same computation with \(\neg B\) gives \(\kappa^{x,y}_A(\neg B)=\kappa(\neg B)\).

\emph{(5)\(\Rightarrow\)(2).}  Let \(N\) be the largest of the four finite
conditional ranks \(\kappa(B\mid A)\), \(\kappa(\neg B\mid A)\),
\(\kappa(B\mid\neg A)\), \(\kappa(\neg B\mid\neg A)\).  The pair
\((x,y)=(0,N)\) is admissible, and
\(\kappa^{0,N}_A(B)=\min\{\kappa(B\mid A),\kappa(B\mid\neg A)+N\}=\kappa(B\mid A)\),
because \(\kappa(B\mid\neg A)+N\geq N\geq\kappa(B\mid A)\).  By~(5),
\(\kappa(B\mid A)=\kappa(B)\), and the same argument with \(\neg B\) gives
\(\kappa(\neg B\mid A)=\kappa(\neg B)\).  The pair \((N,0)\) gives in the
same way \(\kappa(B\mid\neg A)=\kappa(B)\) and
\(\kappa(\neg B\mid\neg A)=\kappa(\neg B)\).  Hence
\(\tau(B\mid A)=\kappa(\neg B)-\kappa(B)=\tau(B\mid\neg A)\).

\emph{Symmetry.}  For the pair \((B,A)\) the four cells are
\(\kappa(B\land A)=a\), \(\kappa(B\land\neg A)=c\),
\(\kappa(\neg B\land A)=b\), \(\kappa(\neg B\land\neg A)=d\), all finite,
and condition~(3) for that pair reads \(a+d=c+b\), the same equation.  By
(1)\(\Leftrightarrow\)(3), applied to both pairs, \(B\) is independent of
\(A\) exactly when \(A\) is independent of \(B\).

\emph{Invariance of the two conditional two-sided ranks.}  Let
\(\kappa'=\kappa^{x,y}_A\).  Since \(\kappa(A\land B\mid\neg A)
=\kappa(\varnothing)-\kappa(\neg A)=\infty\),
\[
 \kappa'(A\land B)=\min\{\kappa(A\land B\mid A)+x,\ \infty\}=a-\kappa(A)+x,
\]
and in the same way \(\kappa'(A\land\neg B)=b-\kappa(A)+x\),
\(\kappa'(\neg A\land B)=c-\kappa(\neg A)+y\) and
\(\kappa'(\neg A\land\neg B)=d-\kappa(\neg A)+y\).  These four cells are
finite, so, by the first paragraph applied to \(\kappa'\),
\(\tau'(B\mid A)=(b-\kappa(A)+x)-(a-\kappa(A)+x)=b-a=\tau(B\mid A)\) and
\(\tau'(B\mid\neg A)=d-c=\tau(B\mid\neg A)\).  The reason relation between
\(A\) and \(B\) depends only on these two values, so it is unchanged.
\end{proof}
\end{annexitem}

\begin{annexitem}{annex:p15-w02b-p15-039-consonant}{thm:p15-w02b-p15-039-consonant}{Consonant ranks and Dempster's rule}
\begin{proof}
(i) Since \(k_j<k_{j+1}\), each \(2^{-k_j}-2^{-k_{j+1}}\) is positive, and
\(2^{-k_t}>0\).  The total mass telescopes:
\(\sum_{j<t}(2^{-k_j}-2^{-k_{j+1}})+2^{-k_t}=2^{-k_0}=1\).  Because
\(\kappa(W)=0\) and \(W\) is finite, Theorem~\ref{thm:p15-rat53-rank-laws}
gives a world of rank \(0=k_0\), so \(F_0\ne\varnothing\) and the empty set
has mass zero.  A world of rank \(k_{j+1}\) lies in \(F_{j+1}\) but not in
\(F_j\), so the inclusions are strict.  For any \(A\),
\(F_j\land A\ne\varnothing\) holds exactly when some \(w\in A\) has
\(r(w)\leq k_j\), that is, exactly when \(\kappa(A)\leq k_j\), since
\(\kappa(A)=\min_{w\in A}r(w)\).  If \(\kappa(A)=\infty\) no index
qualifies and \(\operatorname{Pl}(A)=0=2^{-\infty}\).  Otherwise
\(\kappa(A)\) is a finite value of \(r\), say \(k_i\), the qualifying
indices are \(i,\ldots,t\), and the sum of their masses telescopes to
\(2^{-k_i}\).  This is \(\Pi(A)\) by the definition in
Theorem~\ref{thm:p15-rat53-possibility}.  For belief,
\(F_j\not\subseteq A\) exactly when \(F_j\land\neg A\ne\varnothing\), so
\[
\begin{aligned}
 \operatorname{Bel}(A)&=\sum_{F_j\subseteq A}m_\kappa(F_j)
 =1-\sum_{F_j\land\neg A\ne\varnothing}m_\kappa(F_j)\\
 &=1-\operatorname{Pl}(\neg A)=1-2^{-\kappa(\neg A)}=N(A).
\end{aligned}
\]

(ii) With \(m^E(E)=1\), the conflict-free weight of
Definition~\ref{def:p15-r47-combination} is
\(K=\sum_{F_j\land E\ne\varnothing}m_\kappa(F_j)=\operatorname{Pl}(E)=2^{-\kappa(E)}\),
positive because \(\kappa(E)<\infty\).  The combination gives each nonempty
\(G\) the mass \(K^{-1}\sum_{j:F_j\land E=G}m_\kappa(F_j)\), so its focal
sets are among the sets \(F_j\land E\), which form a chain because the
\(F_j\) do.  For any \(A\), grouping the indices by the set
\(G=F_j\land E\),
\[
\begin{aligned}
 \operatorname{Pl}_{m_\kappa\otimes m^E}(A)
 &=K^{-1}\sum_{j:F_j\land E\land A\ne\varnothing}m_\kappa(F_j)\\
 &=\frac{\operatorname{Pl}(A\land E)}{\operatorname{Pl}(E)}
 =2^{-(\kappa(A\land E)-\kappa(E))}.
\end{aligned}
\]
If \(\kappa(A\land E)<\infty\), the exponent is \(-\kappa(A\mid E)\) by
Definition~\ref{def:p15-rat53-conditionals}.  If
\(\kappa(A\land E)=\infty\), both \(\operatorname{Pl}(A\land E)\) and
\(2^{-\kappa(A\mid E)}\) are zero.

(iii) For \(\lambda_1\) the finite values are \(0,1\), with
\(F_0=\{x,y\}\) and \(F_1=W\), so \(m_{\lambda_1}\) gives \(1/2\) to each.
For \(\lambda_2\), \(F_0=\{y,z\}\) and \(F_1=W\), again with masses
\(1/2\).  Every pair of focal sets meets, because
\(\{x,y\}\land\{y,z\}=\{y\}\) and the other pairs contain \(W\); hence
\(K=1\), and the four products, each \(1/4\), go to \(\{y\}\),
\(\{x,y\}\), \(\{y,z\}\) and \(W\).  Summing the masses of subsets gives
\(\operatorname{Bel}(\{x,y\})=1/4+1/4\),
\(\operatorname{Bel}(\{y,z\})=1/4+1/4\) and
\(\operatorname{Bel}(\{y\})=1/4\).  Suppose
\(m_{\lambda_1}\otimes m_{\lambda_2}=m_\mu\) for a negative ranking
function \(\mu\).  By~(i) its belief is
\(A\mapsto1-2^{-\mu(\neg A)}\).  Since
\(\mu(\neg A\lor\neg B)=\min\{\mu(\neg A),\mu(\neg B)\}\) and
\(s\mapsto1-2^{-s}\) is increasing on \(\mathbb N\cup\{\infty\}\), this belief
satisfies \(\operatorname{Bel}(A\land B)=\min\{\operatorname{Bel}(A),\operatorname{Bel}(B)\}\).
For \(A=\{x,y\}\) and \(B=\{y,z\}\) the left side is \(1/4\) and the right
side is \(1/2\), a contradiction.  The failure agrees with
Theorem~\ref{thm:p15-r47-focal}: the focal sets \(\{x,y\}\) and
\(\{y,z\}\) are not nested.
\end{proof}
\end{annexitem}
